\documentclass[11pt]{article}
\usepackage[margin=1in]{geometry}
\usepackage{amsmath, amssymb, amsthm}
\usepackage{algorithm}
\usepackage{algpseudocode}
\usepackage{enumitem}
\usepackage{hyperref}
\usepackage{bm}

\newtheorem{assumption}{Assumption}
\newtheorem{definition}{Definition}
\newtheorem{proposition}{Proposition}

\DeclareMathOperator*{\argmax}{arg\,max}
\DeclareMathOperator{\LSE}{LogSumExp}       % i don't like you using this as an operator please just show it in math
\DeclareMathOperator{\Vol}{Vol}
\DeclareMathOperator{\CVaR}{CVaR}

\title{Bayesian Inference of Multi-Objective Preferences from Demonstration\\}
\author{Mason duBoef}
\date{April 2026}

\begin{document}
\maketitle

% ============================================================
\section{Problem Setting}
% ============================================================

We consider a multi-objective Markov decision process (MOMDP) with $d$ objectives. At each timestep the environment emits a vector-valued reward $\mathbf{r}(s,a) \in \mathbb{R}^d$. A decision-maker's preferences over the $d$ objectives are encoded by a weight vector $\mathbf{w}$.

We assume access to a precomputed \emph{convex coverage set} (CCS) of Pareto optimal policies. Each policy in the CCS is a stochastic softmax policy trained under a specific preference weight. That being said, each policy is likely optimal for a wide range of preferences. We also have a demonstrator with an unknown preferences $\mathbf{w}_h$. Our goal is to infer $\mathbf{w}_h$ from demonstration data and assign the demonstrator a policy from the CCS that is optimal (or near optimal) under their hidden preferences. We infer a distribution over all possible preference weights. Our approach does not call upon precomputed Q-values or require any RL training to be performed during inference. This ultimately enables us to assign the demonstrator a the pareto optimal policy that corresponds to their inferred preference weight, one which hopefully optimizes their utility under their unknown latent set of preferences.

% ============================================================
\section{Notation}
% ============================================================

\begin{definition}[Multi-Objective MDP]
A multi-objective MDP is a tuple $\langle \mathcal{S}, \mathcal{A}, \mathcal{P}, \bm{\mathcal{R}}, \gamma \rangle$ where $\mathcal{S}$ is the state space, $\mathcal{A}$ is the action space, $\mathcal{P}: \mathcal{S} \times \mathcal{A} \times \mathcal{S} \to [0,1]$ is the transition function, $\bm{\mathcal{R}}: \mathcal{S} \times \mathcal{A} \to \mathbb{R}^d$ is the vector-valued reward function, and $\gamma \in [0,1]$ is the discount factor.
\end{definition}

\begin{definition}[Vector-Valued Q-Function]
For a policy $\pi$, the vector-valued Q-function returns the expected future returns associated with a given state action pair. Since this is a multi objective setting these Q functions return a vector where each entry corresponds to a given objective. They are defined as:
\begin{equation}
    Q^{\pi}(s, a) = \mathbb{E}_{\pi}\!\left[\sum_{t=0}^{\infty} \gamma^t \, \mathbf{r}(s_t, a_t) \;\middle|\; s_0 = s,\; a_0 = a\right] \in \mathbb{R}^d
\end{equation}
\end{definition}

\begin{definition}[Utility]
The utility of taking action $a$ in state $s$ under policy $\pi$ and preference weight $\mathbf{w}$ is the scalarization of the vector-valued Q-function:
\begin{equation}
    u_{\mathbf{w}}^{\pi}(s, a) = \mathbf{w}^\top Q^{\pi}(s, a)
\end{equation}
This is a scalar value representing how good action $a$ is from state $s$ under the set of preferences $w$.
\end{definition}

% does this get used anywhere in our algorithm?
\begin{definition}[Return Vector]
The expected final return of policy $\pi$ beginning from the initial state distribution is:
\begin{equation}
    J^{\pi} = \mathbb{E}_{\pi}\!\left[\sum_{t=0}^{\infty} \gamma^t \, \mathbf{r}(s_t, a_t)\right] \in \mathbb{R}^d
\end{equation}
\end{definition}

Additional notation used throughout:

\begin{center}
\renewcommand{\arraystretch}{1.3}
\begin{tabular}{cl}
    \hline
    \textbf{Symbol} & \textbf{Definition} \\
    \hline
    $\Pi = \{\pi_1, \ldots, \pi_K\}$ & Convex coverage set of $K$ Pareto optimal softmax policies \\
    $\mathbf{w}_k$ & Training weight used to train policy $\pi_k$ \\
    $\beta_{\text{train}}$ & Boltzmann rationality coefficient used in the Pareto optimal policies in the CCS\\
    $\beta_{\text{demo}}$ & Boltzmann rationality coefficient of the demonstrator \\
    $\rho = \beta_{\text{demo}} / \beta_{\text{train}}$ & Rationality ratio \\
    $\mathbf{w}_h$ & Demonstrator's unknown true preference weight \\
    % this could be one trajectory, multiple trajectories or just smaller sequences as actions it doesn't really matter
    $\tau = \{(s_1, a_1), \ldots, (s_T, a_T)\}$ & Observed demonstration data \\
    % so this is the range/space of weight values for which pi k is optimal?
    $R_k \subseteq \Delta^{d-1}$ & Optimality region of policy $\pi_k$ on the weight simplex \\
    % i imagine we have the boundries/range of weights that each policy corresponds to so I imagine getting the volume from that is doable
    $V_k = \Vol(R_k)$ & Volume (Lebesgue measure) of region $R_k$ on the simplex \\
    $\mathbf{c}_k$ & Centroid of region $R_k$ \\
    \hline
\end{tabular}
\end{center}

% ============================================================
\section{Assumptions}
% ============================================================

\begin{assumption}[Fixed Preferences]
\label{a:fixed}
The demonstrator has a fixed but unknown preference $\mathbf{w}_h \in \Delta^{d-1}$ that does not change during the demonstration.
\end{assumption}

\begin{assumption}[Boltzmann-Rational Demonstrator]
\label{a:boltz}
The demonstrator selects actions according to a Boltzmann (softmax) distribution over scalarized Q-values:
\begin{equation}
    P_{\text{demo}}(a \mid s, \mathbf{w}_h) = \frac{\exp\!\big(\beta_{\text{demo}} \cdot \mathbf{w}_h^\top Q^{\pi^*_{\mathbf{w}_h}}(s, a)\big)}{\displaystyle\sum_{a' \in \mathcal{A}(s)} \exp\!\big(\beta_{\text{demo}} \cdot \mathbf{w}_h^\top Q^{\pi^*_{\mathbf{w}_h}}(s, a')\big)}
\end{equation}
where $\pi^*_{\mathbf{w}_h}$ is the Pareto optimal policy corresponding to $\mathbf{w}_h$ and $\beta_{\text{demo}}$ determines the rationality of the demonstrator.
\end{assumption}

\begin{assumption}[Softmax Policy Training]
\label{a:softmax}
% why is the rationality coefficient needed here? Shouldn't these be by definition optimal. I know why we need to model the deonstrator as sub-optimal/boltzman rational but idk why we have to define all the policies in the CCS as boltzman rational too. It seems to me that beta should be near infinity though I understand we need non-zero probability for all actions.
Each policy $\pi_k$ in the CCS was trained as a softmax policy with a set rationality coefficient $\beta_{\text{train}}$, so that for all states $s$ and actions $a$:
\begin{equation}
\label{eq:softmax_policy}
    \pi_k(a \mid s) = \frac{\exp\!\big(\beta_{\text{train}} \cdot \mathbf{w}_k^\top Q^{\pi_k}(s, a)\big)}{\displaystyle\sum_{a' \in \mathcal{A}(s)} \exp\!\big(\beta_{\text{train}} \cdot \mathbf{w}_k^\top Q^{\pi_k}(s, a')\big)}
\end{equation}
This causes $\pi_k(a \mid s) > 0$ for all $s, a$, so all actions have non-zero probability.
\end{assumption}

% do we need to state this explicitly? do we call upon it later?
\begin{assumption}[Conditional Independence]
\label{a:indep}
Actions at different timesteps are conditionally independent given $\mathbf{w}_h$:
\begin{equation}
    P(\tau \mid \mathbf{w}_h) = \prod_{t=1}^{T} P(a_t \mid s_t, \mathbf{w}_h)
\end{equation}
\end{assumption}

% ============================================================
\section{Key Identity: Difference In Utility from Policy Probabilities}
% ============================================================

% show how this identity is derived step by step
From Assumption~\ref{a:softmax}, the log-ratio of action probabilities under policy $\pi_k$ encodes the scalarized Q-value difference (aka. difference in utility):
\begin{equation}
\label{eq:key_identity}
    \mathbf{w}_k^\top Q^{\pi_k}(s, a) - \mathbf{w}_k^\top Q^{\pi_k}(s, a') = \frac{1}{\beta_{\text{train}}} \Big[\log \pi_k(a \mid s) - \log \pi_k(a' \mid s)\Big]
\end{equation}

% we dont need the stored J values here in order to infer this difference in utility between actions at each state
This identity allows us to express the demonstrator's action likelihood entirely in terms of the stored policy probabilities, without accessing Q-values directly.

% when we are calulating differences we only need to calculate |A|-1 differences. We do't need the difference between each pair of 2 actions 

\begin{proposition}[Likelihood in Terms of Policy Probabilities]
\label{prop:likelihood}
Under Assumptions~\ref{a:boltz}--\ref{a:softmax}, for a candidate weight $\mathbf{w}_n$ whose corresponding optimal policy is $\pi_n^*$, the demonstrator's probability of choosing action $a$ in state $s$ is:
\begin{equation}
\label{eq:likelihood}
    P(a \mid s, \mathbf{w}_n) = \frac{\pi_n^*(a \mid s)^{\,\rho}}{\displaystyle\sum_{a' \in \mathcal{A}(s)} \pi_n^*(a' \mid s)^{\,\rho}}
\end{equation}
where $\rho = \beta_{\text{demo}} / \beta_{\text{train}}$ is the rationality ratio.
\end{proposition}

\begin{proof}
Substituting the encoded Q-values from Eq.~\eqref{eq:key_identity} into the Boltzmann likelihood from Assumption~\ref{a:boltz}:
\begin{align}
    P(a \mid s, \mathbf{w}_n) &\propto \exp\!\big(\beta_{\text{demo}} \cdot \mathbf{w}_n^\top Q^{\pi_n^*}(s, a)\big) \nonumber \\
    &= \exp\!\left(\frac{\beta_{\text{demo}}}{\beta_{\text{train}}} \cdot \log \pi_n^*(a \mid s) + c(s)\right) \nonumber \\
    &\propto \exp\!\big(\rho \cdot \log \pi_n^*(a \mid s)\big) = \pi_n^*(a \mid s)^{\,\rho}
\end{align}
where $c(s)$ is a state-dependent constant (from the log-partition function of $\pi_n^*$) that cancels in the normalization. The three behaviorally significant cases of $\rho$ are:
\begin{itemize}[itemsep=2pt]
    \item $\rho = 1$: The demonstrator is as rational as the trained policies. The likelihood reduces to $P(a \mid s, \mathbf{w}_n) = \pi_n^*(a \mid s)$.
    \item $\rho > 1$: The demonstrator is more rational (probabilities sharpen toward the greedy action).
    \item $\rho < 1$: The demonstrator is noisier (probabilities flatten toward uniform).
\end{itemize}
\end{proof}

% ============================================================
\section{Region-Based Posterior}
% ============================================================

Since the likelihood (Eq.~\ref{eq:likelihood}) depends on $\mathbf{w}_n$ only through the identity of the corresponding optimal policy $\pi_n^*$, all candidate weights mapping to the same policy receive identical likelihood updates. This induces a natural partition of the weight simplex.

\begin{definition}[Optimality Regions]
For each policy $\pi_k \in \Pi$, its optimality region is:
\begin{equation}
    R_k = \left\{\mathbf{w} \in \Delta^{d-1} : \mathbf{w}^\top J^{\pi_k} \geq \mathbf{w}^\top J^{\pi_j} \;\;\forall\; j \neq k\right\}
\end{equation}
Each $R_k$ is a convex polytope on the simplex, and $\{R_1, \ldots, R_K\}$ partitions $\Delta^{d-1}$ (up to measure-zero boundaries). The region boundaries are hyperplanes defined by $\mathbf{w}^\top(J^{\pi_k} - J^{\pi_j}) = 0$.

For $d = 2$ objectives, the simplex is the interval $[0,1]$ and each $R_k$ is a sub-interval with analytically computable endpoints, volumes, and centroids.
\end{definition}

Rather than maintaining a posterior over $N$ discretized weight points, we maintain it over $K$ regions. The posterior is piecewise-uniform within each region:
% this equation confuses me I understand that we are looking updating a probability distribution over all weights be dividing the space into regions with 1 centroid representing the region and then dividing the pobroability assigned to that centroid/region uniformly over the weights in that region. But what exactly is this eqation saying, in particualrly what is the fraction at the end representing?
% Also how are the centroid weights calculated? If the region has a weird shape like a donut shape is it possible that the centroid could actually fall inside a different region?
\begin{equation}
    P(\mathbf{w} \mid \tau) = \sum_{k=1}^{K} P(R_k \mid \tau) \cdot \frac{\mathbf{1}[\mathbf{w} \in R_k]}{V_k}
\end{equation}

% ============================================================
\section{Algorithm}
% ============================================================
 
\begin{algorithm}[H]
\caption{BMOPID: Bayesian Multi-Objective Preference Inference from Demonstration}
\label{alg:bmopid}
\begin{algorithmic}[1]
\Require CCS of softmax policies $\Pi = \{\pi_1, \ldots, \pi_K\}$ with return vectors $J^{\pi_k} \in \mathbb{R}^d$
\Require Region volumes $V_k$ and centroids $\mathbf{c}_k$ (precomputed from $J$ vectors)
\Require Demonstration $\tau = \{(s_1, a_1), \ldots, (s_T, a_T)\}$
\Require Rationality ratio $\rho = \beta_{\text{demo}} / \beta_{\text{train}}$
\Ensure Posterior distribution $P(R_k \mid \tau)$ for $k = 1, \ldots, K$
 
\Statex
\Statex \textbf{// Phase 1: Initialize uniform prior}
\For{$k = 1, \ldots, K$}
    \State $\ell_k \gets \log V_k$
        \Comment{larger regions get more prior mass}
\EndFor
 
\Statex
\Statex \textbf{// Phase 2: Sequential Bayesian update}
\For{each $(s_t, a_t) \in \tau$}
    \For{$k = 1, \ldots, K$}
        \State Compute log-likelihood of observed action under policy $\pi_k$:
        \begin{equation*}
            \lambda_{k,t} = \rho \cdot \log \pi_k(a_t \mid s_t) - \log\!\left(\sum_{a' \in \mathcal{A}(s_t)} \pi_k(a' \mid s_t)^{\,\rho}\right)
        \end{equation*}
        \State $\ell_k \gets \ell_k + \lambda_{k,t}$
            \Comment{accumulate log-posterior}
    \EndFor
\EndFor
 
\Statex
% We have the probabilites over all centroids where each centroid represents a region so now we just distrubute the probability mass for each centroid uniformly over the region it represents, bringing us to a continous (but not smooth) distribution
\Statex \textbf{// Phase 3: Normalize}
\State Compute $Z = \log\!\left(\sum_{j=1}^{K} \exp(\ell_j)\right)$
\For{$k = 1, \ldots, K$}
    \State $P(R_k \mid \tau) \gets \exp(\ell_k - Z)$
\EndFor
 
\Statex
\Statex \textbf{// Phase 4: Extract outputs}
\State $\hat{k} \gets \argmax_k P(R_k \mid \tau)$
    \Comment{MAP region}
\State $\hat{\mathbf{w}}_{\text{MAP}} \gets \mathbf{c}_{\hat{k}}$
    \Comment{MAP weight estimate (centroid of highest-probability region)}
\State $\bar{\mathbf{w}} \gets \sum_{k=1}^{K} P(R_k \mid \tau) \cdot \mathbf{c}_k$
    \Comment{posterior mean weight (optional)}
\end{algorithmic}
\end{algorithm}


\newpage

% ============================================================
\section{Policy Selection}
\label{sec:policy_selection}
% ============================================================

% I see 4 reasonable ways to assign a policy given our posterior
    % we could call the region with the highest probability the "inferred weight" and assign the corresponding policy
    % we could calculate the mean weight (weighted by the probabilites in the posterior) call that the "inferred weight" and assign the corresponding policy
    % we could assign the policy that gives the highest expected utility given the whole probability distribion over weights
    % we could do something risk-sensitive and assign the policy with the highest 5% CVaR utility

Given the posterior $P(R_k \mid \tau)$, there are three natural ways to assign a policy to the demonstrator.
 
\subsection{MAP Policy}
 
The simplest approach: assign the policy corresponding to the highest-probability region.
\begin{equation}
    \hat{\pi}_{\text{MAP}} = \pi_{\hat{k}}, \quad \hat{k} = \argmax_k P(R_k \mid \tau)
\end{equation}
This is appropriate when the posterior is sharply peaked and we are confident in the inferred preference.
 
\subsection{Maximum Expected Utility Policy}
 
Assign the policy that maximizes expected utility under the full posterior. For each candidate policy $\pi_k$, its expected utility across the posterior is:
\begin{equation}
    U(\pi_k) = \sum_{j=1}^{K} P(R_j \mid \tau) \cdot \mathbf{c}_j^\top J^{\pi_k}
\end{equation}
The selected policy is:
\begin{equation}
    \hat{\pi}_{\text{EU}} = \argmax_{\pi_k \in \Pi} \sum_{j=1}^{K} P(R_j \mid \tau) \cdot \mathbf{c}_j^\top J^{\pi_k}
\end{equation}
This accounts for the full shape of the posterior rather than just its mode, and can select a different policy than MAP when the posterior is spread across multiple regions.
 
\subsection{Risk-Sensitive Policy (CVaR)}
 
A risk-sensitive approach that guards against the worst plausible outcomes. CVaR at level $\alpha$ (e.g., $\alpha = 0.05$) focuses on the worst $\alpha$-fraction of the posterior and asks: what is the expected utility in that worst-case tail?
 
For each candidate policy $\pi_k$, compute the utility $\mathbf{w}^\top J^{\pi_k}$ at each region centroid, weighted by posterior mass. Sort these by utility in ascending order. The CVaR is the probability-weighted average utility across the worst regions whose cumulative posterior mass sums to $\alpha$:
\begin{equation}
    \text{CVaR}_\alpha(\pi_k) = \frac{1}{\alpha} \sum_{j \in T_\alpha(\pi_k)} P(R_j \mid \tau) \cdot \mathbf{c}_j^\top J^{\pi_k}
\end{equation}
where $T_\alpha(\pi_k)$ is the set of regions (sorted by ascending $\mathbf{c}_j^\top J^{\pi_k}$) whose cumulative posterior mass first reaches $\alpha$. The selected policy is:
\begin{equation}
    \hat{\pi}_{\text{CVaR}} = \argmax_{\pi_k \in \Pi}\; \text{CVaR}_\alpha(\pi_k)
\end{equation}
This selects the policy with the best worst-case expected utility, hedging against the possibility that the demonstrator's true preference lies in a region where this policy performs poorly.


\end{document}
